\begin{frame}{Formel für Darstellungslänge in Basis $b$}
	Es seien $n>0$ und $b>1$ natürliche Zahlen. Die kürzeste $b$-adische Darstellung von $n$ (Darstellung ohne führende Nullen) hat genau:
	\begin{align*}
		\myhighlight{Goldenrod}{$\ceil{\log_b(n+1)}$}
	\end{align*}
	Stellen.

	\vspace{0.5cm}
	Zur Erinnerung: Für eine reelle Zahl $x$, ist $\ceil{x}$ die kleinste ganze Zahl $\geq x$ (\enquote{Aufrundeklammer}).
\end{frame}

\begin{frame}{Rechenaufgabe}
	\begin{myexercise}\label{myexercise:001}
		Wie viele Stellen hat die kürzeste Darstellung der Zahl $83$ in Basis $b = 3$?
	\end{myexercise}
\end{frame}

\begin{frame}{Lösung zur Rechenaufgabe}
	\begin{myanswer}
		Die Anzahl $m$ der Stellen in der kürzesten Darstellung der Zahl $83$ in Basis $3$ ist
		\begin{align*}
			m = \ceil{\log_3(83 + 1)}.
		\end{align*}
		Da $3^4 = 81$ und $3^5 = 243$ folgt, dass $4 < \log_3(83 + 1) < 5$ und somit
		\begin{align*}
			m = \ceil{\log_3(83 + 1)} = 5.
		\end{align*}
	\end{myanswer}
\end{frame}

\begin{frame}{Wissenssicherung}
	\begin{myexercise}\label{myexercise:002}
		Wie viele Stellen hat die kürzeste Darstellung der Zahl $500$ in Basis $b = 16$?
	\end{myexercise}
\end{frame}

\begin{frame}{Lösung zur Wissenssicherung}
	\begin{myanswer}
		Die Anzahl $m$ der Stellen in der kürzesten Darstellung der Zahl $500$ in Basis $16$ ist
		\begin{align*}
			m = \ceil{\log_{16}(500 + 1)}.
		\end{align*}
		Da $16^2 = 256$ und $16^3 = 4096$ folgt, dass $2 < \log_{16}(500 + 1) < 3$ und somit
		\begin{align*}
			m = \ceil{\log_{16}(500 + 1)} = 3.
		\end{align*}
	\end{myanswer}
\end{frame}

\begin{frame}{Wissenssicherung}
	\begin{myexercise}\label{myexercise:003}
		Wie viele Stellen hat die kürzeste Darstellung der Zahl $8^5 - 1$ in Basis $b = 2$?
	\end{myexercise}
\end{frame}

\begin{frame}{Lösung zur Wissenssicherung}
	\begin{myanswer}
		Die Anzahl $m$ der Stellen in der kürzesten Darstellung der Zahl $8^5 - 1$ in Basis $2$ ist
		\begin{align*}
			m = \ceil{\log_{2}\lr{8^5}} = \ceil{\log_2\lr{2^{15}}} = 15.
		\end{align*}
	\end{myanswer}
\end{frame}


